\caption{Matched clinical-only and paired cascades at calibration-selected workload}\label{tab:matched}\centering
\small\renewcommand{\arraystretch}{1.00}\setlength{\tabcolsep}{3pt}
\begin{tabular*}{\textwidth}{@{\extracolsep{\fill}}llrrrrrrr@{}}\toprule
Cohort & Head & Images & \multicolumn{2}{c}{TP} & \multicolumn{2}{c}{FP} & \multicolumn{2}{c}{Floors met}\\
& & & Paired & Clinical & Paired & Clinical & Paired & Clinical\\\midrule
MRA & Score & 154.9 & 372.8 & 373.0 & 348.8 & 351.0 & 0/9 & 0/9\\
MRA & Feature & 159.7 & 369.3 & 366.0 & 353.3 & 356.0 & 0/9 & 0/9\\
HIBA & Feature & 50.0 & 184.9 & 185.0 & 64.1 & 64.3 & 7/9 & 6/9\\
UQ & Score & 1352.2 & 90.8 & 90.3 & 6304.9 & 6305.7 & 0/9 & 0/9\\
UQ & Feature & 1335.4 & 89.8 & 89.0 & 6205.1 & 6208.0 & 0/9 & 0/9\\
MCR-SL & Score & 38.2 & 35.7 & 36.2 & 106.8 & 106.8 & 0/9 & 0/9\\
MCR-SL & Feature & 40.0 & 37.4 & 37.2 & 96.6 & 96.8 & 3/9 & 3/9\\
\bottomrule\end{tabular*}\par\smallskip\parbox{\linewidth}{\footnotesize Counts average nine systems. Clinical replay uses identical routing/flags and no evaluation dermoscopy; both arms use paired development/calibration. Score/Feature controls flag the top $L$ clinical cases (ID ties). For frozen-feature controls, design validation selects PCA32/full features and linear/RBF SVMs ($C=1$) by top-half melanoma recall, then AUC; heads balance classes/cohorts. HIBA Score appears in Table~\ref{tab:localworkload}A. Heads need not share workload. The complete subtype/system/fold counts and paired group-bootstrap intervals remain in the data archive.}

